\section{Introduction}

Modern data-intensive services increasingly rely on Approximate Nearest Neighbor (ANN) search over high-dimensional embeddings. However, scaling ANN search to high recall under strict latency SLOs remains fundamentally constrained by memory capacity and I/O bandwidth.
This tension is now a dominant systems bottleneck: improving recall predictably inflates latency and operational cost, while controlling latency requires either aggressive approximation or prohibitively large memory footprints. Consequently, ANN search has become a critical performance and cost bottleneck in production systems, underpinning recommendation pipelines~\cite{Covington:2016:DNN:YouTube}, search engines~\cite{Liu:2021:Q2S:Facebook, amazon-shop-the-look, visrel}, and Retrieval-Augmented Generation (RAG) workloads for large language models (LLMs)~\cite{Lewis:2020:RAG:NeurIPS, gao2024retrievalaugmentedgenerationlargelanguage}, more recently including agentic applications~\cite{zhao2026amabench, lee2026characterizing}.

This bottleneck is structural rather than incidental. Achieving high recall in ANN search requires traversing increasingly large and irregular neighborhoods in high-dimensional space, leading to random access amplification that cannot be efficiently prefetched or batched \cite{aker}. In graph-based indexes, recall improvements translate directly into expanded graph exploration, candidate explosion, and deeper disk or memory accesses. Although in-memory ANN systems achieve low latency, their memory cost grows linearly with the size of the dataset. Disk-based systems reduce memory pressure, but suffer severe latency degradation at high recall, especially on very large datasets such as billion-scale collections, where pointer chasing and random I/O \cite{jaiswal2022ooddiskannefficientscalablegraph} dominate execution time.  As a result, operators are forced into an undesirable choice between cost and performance, with no mechanism to amortize repeated query work over time. 

Applications that rely on vector search for retrieval exhibit substantial locality. However, unlike the temporal locality observed in traditional database workloads~\cite{chandramouli2018faster}, where the same data objects are repeatedly accessed over time, locality in vector search does not primarily arise from exact query repetition. Instead, repetition occurs in the latent space, where semantically related queries map to nearby embedding regions and retrieve overlapping sets of vectors. We refer to this behavior as \emph{temporal-semantic locality}. A recent study~\cite{lee2026characterizing} reports that vector search queries target overlapping regions of the embedding space and repeatedly retrieve the same vectors, with 73\% of the data vectors accessed by coding agents being accessed repeatedly within and across sessions. This phenomenon also occurs in traditional workloads: up to 40\% and 68\% of user queries in web search and recommender systems, respectively, are semantically repetitive variants of previous queries and thus have nearby vector embeddings~\cite{10.1145/1277741.1277770yahoo, gill2025advancingsemanticcachingllms, li2024recommenderpurposerepeatexploration}.

Systems exploit this locality to reduce query latency through caching. One approach, \emph{traversal-centric caching}, keeps auxiliary index structures in memory, such as data pages, IVF clusters, or graph nodes near the entry point~\cite{pgvector, diskann, 3495724.3496619, 10.1145/3318464.3386131}. Orthogonal to this approach, \emph{result-centric caching} leverages approximate, or similarity, caching~\cite{similarity-caching, similarity-caching-theory-algorithms}. Such systems ~\cite{aker} serve queries directly from cached results of previous queries when these results can answer a new query with sufficient confidence, typically determined by a distance threshold. By bypassing the underlying index traversal, result-centric caching can reduce cache-hit latency by orders of magnitude while introducing negligible memory overhead at the caching layer.

